\subsection{Exponentiation as an order on finitely supported functions}
\label{sec:exporder}
Fix canonical terms $a\ge2$ and $b$. A function from $I_b$ to $I_a$
of finite support is represented by a list of pairs $(s,d)$ in strictly
decreasing order of $s<b$, with $0<d<a$. At every unlisted argument its
value is zero. Order two distinct lists by their digits at the largest
argument where they differ. That argument is found by merging their
two finite lists of occupied positions.

\begin{theorem}[The order-type interpretation of ordinal powers]
Uniformly in $a,b$, the map
\[
                f\longmapsto V(f)=\sum_{s\text{ in decreasing order}}
                                      a^s f(s)
\]
is a primitive recursive order isomorphism onto $I_{a^b}$, with a
primitive recursive inverse. This includes the empty support and $b=0$.
\end{theorem}
\begin{proof}
First allow arbitrary finite lists of occupied positions, without an
upper bound $b$. For a nonempty list with highest position $s$ its
value lies in $[a^s,a^{s+1})$. To prove this, consider its first
digit $d$ and the value $L$ of the tail. If the tail is empty,
$L=0<a^s$. Otherwise its highest position $t<s$ gives
$L<a^{t+1}\le a^s$ by induction on the finite list length. Thus
\[
 a^s\le a^s d\le a^s d+L<a^s d+a^s=a^s(d+1)\le a^{s+1}.
\]
Here $d+1\le a$ follows from the successor property. Likewise the
tail strictly below any position $s$ has value less than $a^s$.
At the greatest differing position write the two digits as $d<d'$,
allowing zero. The smaller digit and its tail then contribute strictly
less than $a^s(d+1)\le a^s d'$. Restoring the identical higher terms
by addition on the left preserves the strict inequality. This proves
order preservation and reflection, and gives injectivity. The same
bounds show that $V(f)<a^b$ exactly when all occupied positions are
below $b$, including the empty list since $a^b>0$.

It remains to produce the digits of any given $x$. Put $k=k_a$ as in
\eqref{eq:arith-power}, so $a^\omega=W(k)$ and
$a^{\omega J+i}=W(kJ)a^i$ for finite $i$. For each outer exponent $e$
of $x$, divide on the left by $k$, writing $e=kJ+r$ with $r<k$.
The intervals $[kJ,k(J+1))$ are ordered by $J$, and inside such an
interval the order is the order of its remainders. Consecutive entries
with a common $J$ therefore give a canonical word $C_J$ of their
remainders, including repeated entries and zeros. Its head is $<k$, so
$0<C_J<W(k)$. Each original group has value $W(kJ)C_J$, since
$W(kJ)W(r)=W(kJ+r)$, including $r=0$. Summing these groups in
decreasing order of $J$ reconstructs $x$.

Find a natural $N$ with $C_J<a^N$ as follows. For finite $a$, $k=1$,
so all entries of $C_J$ are zero. Its value is its length $m$,
and $N=m+1$ gives $a^N>m$. For infinite $a$ put $g=h(h(a))$.
Then $k=W(g+1)$. Every entry of the head $r_*=h(C_J)$ is at most
$g$. With $N=|r_*|+1$, ordinary finite multiplication shows that
$h(a)N$ begins with at least $N$ copies of $g$. Thus
$h(a)N>r_*$, by the first differing exponent or by length when all
of $r_*$ consists of copies of $g$. Since $h(a^N)=h(a)N$,
again $C_J<a^N$.

Now set $q_0=C_J$. For $i<N$ divide $q_i=aq_{i+1}+d_i$,
where $d_i<a$. From $q_i<a^{N-i}$ one gets
$q_{i+1}<a^{N-i-1}$, since $aq_{i+1}\le q_i$ and
$a\,a^{N-i-1}=a^{N-i}$. At the last step $q_N<1$, hence $q_N=0$.
Substituting the equations back, using left distributivity at each
step, gives $C_J=a^{N-1}d_{N-1}+\cdots+a^0d_0$.
Multiplication on the left by $W(kJ)$ moves these digits to the positions
$\omega J+i$. These positions increase with $i$, and positions from
distinct groups are ordered by $J$, since
$\omega J+i<\omega(J+1)$. Omitting zero digits and concatenating the
groups therefore returns a valid finite support. It has value $x$.
The initial-segment bound already proved guarantees the required
positions $<b$ whenever $x<a^b$. Uniqueness follows from the strict
comparison proved above.

There are at most $|x|$ groups, and each group has its displayed
numerical bound $N$ on the division loop. Intermediate lists are
computed by ordinary primitive recursion with a changing accumulator.
This construction uses no induction on the ordinal quotient.
\end{proof}
For $a=1$ there is only the zero-valued function, and $1^b=1$.
For $a=0$ there is one empty function when $b=0$ and none when $b>0$,
in agreement with the exceptional powers. In this last case a total
function on a nonempty domain cannot choose a value at its least input.
